% SNT Seconde - Web, seance 2 - Beamer First Principles
% Compiler avec XeLaTeX sur Overleaf.
\documentclass[aspectratio=169,11pt]{beamer}
\usepackage{fontspec}
\setmainfont{Latin Modern Sans}
\setsansfont{Latin Modern Sans}
\setmonofont{Latin Modern Mono}
\usepackage{polyglossia}
\setmainlanguage{french}
\usepackage{tikz}
\usetikzlibrary{arrows.meta,positioning,calc,shapes.geometric}
\usepackage{listings}
\usepackage{xcolor}
\definecolor{foret}{HTML}{1C322D}
\definecolor{creme}{HTML}{F8F3EE}
\definecolor{or}{HTML}{EBB552}
\definecolor{rose}{HTML}{F1CDBE}
\definecolor{vertclair}{HTML}{DDE9DF}
\definecolor{encre}{HTML}{24372E}
\setbeamercolor{normal text}{fg=creme,bg=foret}
\setbeamercolor{frametitle}{fg=or,bg=foret}
\setbeamercolor{title}{fg=creme}
\setbeamercolor{block title}{fg=foret,bg=or}
\setbeamercolor{block body}{fg=encre,bg=creme}
\setbeamercolor{alerted text}{fg=or}
\setbeamertemplate{navigation symbols}{}
\setbeamertemplate{footline}{\hfill\color{or}\insertframenumber\hspace{0.5cm}\vspace{0.25cm}}
\setbeamertemplate{frametitle}{\vspace{0.35cm}\insertframetitle\par\vspace{0.15cm}\color{or}\rule{\linewidth}{0.6pt}}
\setlength{\parskip}{5pt}
\lstdefinestyle{webcode}{basicstyle=\ttfamily\small,keywordstyle=\color{foret},stringstyle=\color{foret},
  backgroundcolor=\color{creme},frame=single,rulecolor=\color{or},framerule=1pt,
  xleftmargin=5pt,xrightmargin=5pt,columns=fullflexible,keepspaces=true,showstringspaces=false}
\newcommand{\question}[1]{\begin{beamercolorbox}[rounded=true,sep=12pt]{block body}\Large\textbf{#1}\end{beamercolorbox}}
\newcommand{\reponse}[1]{\uncover<2->{\vspace{5mm}\begin{beamercolorbox}[rounded=true,sep=8pt]{block body}\textbf{À retenir :} #1\end{beamercolorbox}}}
\newcommand{\etiquette}[1]{{\color{or}\small\MakeUppercase{#1}}}
\newcommand{\webwindow}[2]{\begin{tikzpicture}[x=1cm,y=1cm]
  \fill[creme,rounded corners=3pt] (0,0) rectangle (10,3.3);
  \fill[vertclair] (0,2.83) rectangle (10,3.3);
  \fill[rose] (.25,3.07) circle (.07);\fill[or] (.49,3.07) circle (.07);
  \node[anchor=west,text=encre,font=\scriptsize\ttfamily] at (.85,3.07) {#1};
  \node[anchor=center,text=foret,align=center,text width=9cm] at (5,1.6) {#2};
  \end{tikzpicture}}
\title{Du texte à la page Web}
\subtitle{SNT Seconde · Le Web · HTML et CSS}
\author{Séance 2 · Expérimenter, observer, comprendre}
\date{}
\begin{document}

\begin{frame}[plain]
  \vfill
  {\color{or}\Large SCIENCES NUMÉRIQUES ET TECHNOLOGIE}\par\vspace{9mm}
  {\Huge\bfseries Du texte à la page Web}\par\vspace{4mm}
  {\Large HTML et CSS, pas à pas}\par\vspace{12mm}
  \begin{tikzpicture}[node distance=5mm]
    \node[draw=or,rounded corners,fill=creme,text=foret,inner sep=10pt] (a) {\texttt{<h1>Bonjour</h1>}};
    \node[right=10mm of a,draw=or,rounded corners,fill=creme,text=foret,inner sep=10pt] (b) {\Large\bfseries Bonjour};
    \draw[-{Stealth},or,thick] (a)--(b);
  \end{tikzpicture}\vfill
\end{frame}

\begin{frame}{La question de départ}
  \etiquette{Problème concret}\par\vspace{5mm}
  \question{Peut-on changer ce qu'affiche un site Web sans changer le véritable site ?}
  \vspace{8mm}
  \begin{center}\webwindow{lasallesaintdenis.com}{\Huge\bfseries LE MOT DES DIRECTEURS}\end{center}
  \vspace{2mm}{\color{rose}Expérience : ouvrir le site du lycée et choisir un titre.}
\end{frame}

\begin{frame}{Expérience : modifier un titre}
  \begin{columns}[T]
  \column{.50\linewidth}
    \begin{block}{Dans le navigateur}
      \begin{enumerate}
      \item Clic droit sur le titre, puis \textbf{Inspecter}.
      \item Cliquer sur la flèche dans le \textbf{carré en pointillés}.
      \item Cliquer sur le titre de la page.
      \item Double-cliquer sur son texte dans le code.
      \item Écrire son prénom, puis Entrée.
      \end{enumerate}
    \end{block}
  \column{.45\linewidth}
    \vspace{1cm}
    \begin{tikzpicture}
      \node[draw=or,fill=creme,text=foret,rounded corners,minimum width=4.5cm,minimum height=1.4cm,align=center] (n) {\footnotesize\texttt{<h2>Ahmed</h2>}\\\Huge\bfseries Ahmed};
      \node[below=5mm of n,text=rose,align=center] {Le titre change\textbf{à l'écran}.};
    \end{tikzpicture}
  \end{columns}
\end{frame}

\begin{frame}{Une prédiction avant F5}
  \question{Si je recharge la page avec F5, le prénom restera-t-il ?}
  \vspace{10mm}
  \begin{enumerate}
    \item Oui, car j'ai changé le code du site.
    \item Non, car j'ai changé seulement mon affichage.
    \item Tous les visiteurs verront mon prénom.
  \end{enumerate}
  \reponse{Réponse 2. F5 reconstruit la page à partir des ressources d'origine ou du cache. La modification locale disparaît.}
\end{frame}

\begin{frame}{Où se trouve la page d'origine ?}
  \begin{center}
  \begin{tikzpicture}[>=Stealth,node distance=1.4cm,
  box/.style={draw=or,fill=creme,text=foret,rounded corners,minimum height=1.6cm,text width=3.2cm,align=center}]
    \node[box] (s) {\Large Serveur Web\\\small fichiers du site};
    \node[box,right=3cm of s] (b) {\Large Navigateur\\\small copie affichée};
    \draw[->,very thick,or] ([yshift=3mm]s.east)--node[above,align=center,text=creme]{ressources\small HTML et CSS} ([yshift=3mm]b.west);
    \draw[<-,thick,rose] ([yshift=-4mm]s.east)--node[below,text=creme]{requête} ([yshift=-4mm]b.west);
    \node[below=1.0cm of b,draw=rose,rounded corners,text=creme,text width=4cm,align=center] {Inspecter modifie la représentation locale de la page.};
  \end{tikzpicture}
  \end{center}
  \pause\begin{block}{Pourquoi est-ce important ?}Une capture d'écran peut montrer un contenu modifié localement.\end{block}
\end{frame}

\begin{frame}{Un nouveau problème : trois phrases}
  \begin{block}{Un fichier texte}
    Bienvenue sur mon site\\[2mm]
    Je suis élève de seconde.\\[2mm]
    Wikipédia
  \end{block}
  \vspace{7mm}
  \question{Comment le navigateur peut-il savoir ce qui est un titre, un paragraphe ou un lien ?}
  \reponse{Il faut décrire le rôle de chaque élément. C'est le principe du HTML.}
\end{frame}

\begin{frame}[fragile]{HTML : donner une structure au texte}
  \begin{columns}[T]
  \column{.52\linewidth}
  \etiquette{Le fichier texte}\par
  \begin{lstlisting}[style=webcode]
<h1>Bienvenue sur mon site</h1>
<p>Je suis eleve de seconde.</p>
  \end{lstlisting}
  \column{.44\linewidth}
  \etiquette{Ce que le navigateur affiche}\par
  \begin{beamercolorbox}[sep=12pt,rounded=true]{block body}
    {\Large\bfseries Bienvenue sur mon site}\par\vspace{5mm}
    Je suis élève de seconde.
  \end{beamercolorbox}
  \end{columns}
  \vspace{8mm}\pause
  \begin{block}{Principe}\texttt{<h1>} annonce un titre principal. \texttt{<p>} annonce un paragraphe.\end{block}
\end{frame}

\begin{frame}{Une balise ouvre et ferme un contenu}
  \begin{center}
  \begin{tikzpicture}[node distance=0mm]
    \node[draw=or,fill=creme,text=foret,rounded corners,minimum height=1.4cm,text width=2.5cm,align=center] (a) {\Large\ttfamily <p>};
    \node[right=0mm of a,draw=or,fill=vertclair,text=foret,minimum height=1.4cm,text width=4.2cm,align=center] (b) {\Large\ttfamily Bonjour};
    \node[right=0mm of b,draw=or,fill=creme,text=foret,rounded corners,minimum height=1.4cm,text width=2.5cm,align=center] (c) {\Large\ttfamily </p>};
    \node[below=4mm of a,text=or] {Ouverture};
    \node[below=4mm of b,text=or] {Contenu};
    \node[below=4mm of c,text=or] {Fermeture};
  \end{tikzpicture}
  \end{center}
  \vspace{14mm}\question{À quoi sert la barre oblique dans \texttt{</p>} ?}
  \reponse{Elle indique que le paragraphe se termine.}
\end{frame}

\begin{frame}[fragile]{Le squelette d'une page HTML}
  \begin{columns}[T]
  \column{.53\linewidth}
  \begin{lstlisting}[style=webcode]
<!DOCTYPE html>
<html lang="fr">
  <head>
    <title>Ma premiere page</title>
  </head>
  <body>
    <h1>Bienvenue</h1>
    <p>Bonjour !</p>
  </body>
</html>
  \end{lstlisting}
  \column{.43\linewidth}
  \begin{block}{\texttt{head}}Informations sur la page. Le texte de \texttt{title} apparaît dans l'onglet.\end{block}
  \begin{block}{\texttt{body}}Contenu que le navigateur affiche dans la fenêtre.\end{block}
  \end{columns}
\end{frame}

\begin{frame}{Six balises à reconnaître}
  \begin{columns}[T]
  \column{.49\linewidth}
    \begin{block}{Contenu}
    \texttt{<h1>} : titre principal\\[2mm]
    \texttt{<p>} : paragraphe\\[2mm]
    \texttt{<strong>} : texte important
    \end{block}
  \column{.49\linewidth}
    \begin{block}{Navigation et média}
    \texttt{<a href="...">} : lien\\[2mm]
    \texttt{<br>} : retour à la ligne\\[2mm]
    \texttt{<img src="..." alt="...">} : image
    \end{block}
  \end{columns}
  \vspace{5mm}\pause
  \begin{block}{À observer}\texttt{br} et \texttt{img} n'entourent pas de contenu. Ils n'ont pas de balise fermante.\end{block}
\end{frame}

\begin{frame}[fragile]{Un lien : du texte et une destination}
  \begin{lstlisting}[style=webcode]
<a href="https://fr.wikipedia.org">Wikipedia</a>
  \end{lstlisting}
  \vspace{7mm}
  \begin{center}\begin{tikzpicture}[>=Stealth]
  \node[draw=or,fill=creme,text=foret,rounded corners,text width=4cm,align=center,minimum height=1.4cm] (t) {\Large Wikipédia\\\small texte cliquable};
  \node[draw=or,fill=creme,text=foret,rounded corners,text width=4cm,align=center,minimum height=1.4cm,right=1.6cm of t] (u) {\ttfamily https://fr.wikipedia.org\\\small destination};
  \draw[->,very thick,or] (t)--node[above,text=creme]{\ttfamily href} (u);
  \end{tikzpicture}\end{center}
  \pause\begin{block}{Deux mots de vocabulaire}\texttt{a} est une balise. \texttt{href} est un attribut qui précise la destination.\end{block}
\end{frame}

\begin{frame}[fragile]{Les trois destinations d'un lien}
  \begin{block}{1. Un autre site}
  \texttt{href="https://fr.wikipedia.org"}
  \end{block}
  \begin{block}{2. Une page du même site}
  \texttt{href="connexion.html"}
  \end{block}
  \begin{block}{3. Un endroit de la page actuelle}
  \texttt{href="\#haut"} vise un élément qui porte \texttt{id="haut"}.
  \end{block}
  \vspace{3mm}\etiquette{Question orale : quelle destination pour contact.html ?}
\end{frame}

\begin{frame}{Changer durablement son propre fichier}
  \begin{center}
  \begin{tikzpicture}[>=Stealth,node distance=14mm,
    bx/.style={draw=or,fill=creme,text=foret,rounded corners,text width=3cm,minimum height=1.5cm,align=center}]
    \node[bx] (e) {\Large Geany / VS Code\\\small modifier \texttt{index.html}};
    \node[bx,right=of e] (f) {\Large Ctrl+S\\\small enregistrer};
    \node[bx,right=of f] (n) {\Large Navigateur\\\small F5 pour afficher};
    \draw[->,thick,or] (e)--(f);\draw[->,thick,or] (f)--(n);
  \end{tikzpicture}
  \end{center}
  \vspace{8mm}
  \begin{block}{Activité 2.3}\Télécharger \texttt{demo-web}, ouvrir \texttt{index.html} dans l'éditeur et dans le navigateur, modifier une phrase, créer un lien, puis afficher \texttt{photo.jpg}.\end{block}
\end{frame}

\begin{frame}{Le problème des 500 pages}
  \question{Un site possède 500 pages avec un titre bleu. Comment rendre tous ces titres rouges rapidement ?}
  \vspace{9mm}
  \begin{center}
  \begin{tikzpicture}[>=Stealth]
  \foreach \i in {0,1,2,3,4}{\node[draw=or,fill=creme,text=foret,rounded corners,minimum width=1.25cm,minimum height=2cm] at (1.5*\i,0) {\ttfamily h1};}
  \node[text=creme] at (8.7,0) {\Large ... 500};
  \end{tikzpicture}
  \end{center}
  \reponse{On écrit une règle de présentation commune dans un fichier CSS.}
\end{frame}

\begin{frame}[fragile]{CSS : une règle pour la présentation}
  \begin{columns}[T]
  \column{.45\linewidth}
  \begin{lstlisting}[style=webcode]
h1 {
  color: red;
}
  \end{lstlisting}
  \column{.52\linewidth}
  \begin{block}{Lire la règle}
    \textbf{À qui ?} \texttt{h1}, les titres\\[2mm]
    \textbf{Quoi ?} \texttt{color}, la couleur du texte\\[2mm]
    \textbf{Comment ?} \texttt{red}, en rouge
  \end{block}
  \end{columns}
  \vspace{8mm}\question{Pour centrer tous les titres, change-t-on le HTML ou le CSS ?}
  \reponse{Le CSS, avec \texttt{text-align: center;}.}
\end{frame}

\begin{frame}{Le même contenu, deux présentations}
  \begin{columns}
  \column{.48\linewidth}
  \begin{beamercolorbox}[rounded=true,sep=12pt]{block body}
  {\color{blue}\Large\bfseries Mon lycée}\par\vspace{4mm}
  {\color{black}Je suis élève de seconde.}
  \end{beamercolorbox}
  \centering\etiquette{Style A}
  \column{.48\linewidth}
  \begin{beamercolorbox}[rounded=true,sep=12pt]{block body}
  {\color{red}\Large\bfseries Mon lycée}\par\vspace{4mm}
  {\color{black}\large Je suis élève de seconde.}
  \end{beamercolorbox}
  \centering\etiquette{Style B}
  \end{columns}
  \vspace{12mm}\pause
  \begin{block}{Séparation des rôles}\HTML décrit ce que contient la page. CSS décrit son apparence.\end{block}
\end{frame}

\begin{frame}{Une feuille de style partagée}
  \begin{center}
  \begin{tikzpicture}[>=Stealth,
   page/.style={draw=or,rounded corners,fill=creme,text=foret,text width=3.0cm,minimum height=1.5cm,align=center}]
   \node[page] (a) at (0,1.3) {\ttfamily index.html};
   \node[page] (b) at (0,-1.3) {\ttfamily connexion.html};
   \node[page,fill=rose] (c) at (6,0) {\ttfamily style.css\\\small règle commune};
   \draw[->,very thick,or] (a.east)--(c.north west);
   \draw[->,very thick,or] (b.east)--(c.south west);
  \end{tikzpicture}
  \end{center}
  \vspace{4mm}
  \begin{block}{Expérience}\Changer \texttt{navy} en \texttt{red} dans \texttt{style.css}, enregistrer, puis recharger les deux pages.\end{block}
\end{frame}

\begin{frame}{Un site lisible sur téléphone ?}
  \question{Faut-il créer une deuxième version de tout le site pour les petits écrans ?}
  \vspace{8mm}
  \begin{center}\begin{tikzpicture}
  \node[draw=or,fill=creme,text=foret,minimum width=5cm,minimum height=2.4cm,align=center] at (0,0) {\Huge\bfseries Titre\\Texte de la page};
  \node[draw=or,fill=creme,text=foret,rounded corners=7pt,minimum width=1.7cm,minimum height=3.2cm,align=center] at (6,0) {\bfseries Titre\\\small Texte};
  \end{tikzpicture}\end{center}
  \reponse{Une seule page HTML peut utiliser des règles CSS adaptées à la largeur de l'écran. C'est le responsive web design, popularisé par Ethan Marcotte en 2010.}
\end{frame}

\begin{frame}{Mini-QCM : qui fait quoi ?}
  \question{Pour chaque demande, choisis HTML ou CSS.}
  \vspace{5mm}
  \begin{enumerate}
  \item Ajouter une photo du lycée.
  \item Mettre tous les liens en vert.
  \item Changer le texte d'un bouton.
  \item Adapter la taille du texte au téléphone.
  \end{enumerate}
  \uncover<2->{\begin{block}{Correction expliquée}1. HTML : ajouter du contenu. \quad 2. CSS : modifier la couleur.\\3. HTML : changer les mots. \quad 4. CSS : adapter la présentation.\end{block}}
\end{frame}

\begin{frame}{La carte mentale de la séance}
  \begin{center}
  \begin{tikzpicture}[>=Stealth,node distance=8mm,
   n/.style={draw=or,fill=creme,text=foret,rounded corners,align=center,text width=3cm,minimum height=1.35cm}]
   \node[n] (f) {\Large Fichiers du site};
   \node[n,below left=1cm and 1.7cm of f] (h) {\Large HTML\\\small structure et contenu};
   \node[n,below right=1cm and 1.7cm of f] (c) {\Large CSS\\\small présentation};
   \node[n,below=2.5cm of f] (b) {\Large Navigateur\\\small page affichée};
   \draw[->,thick,or] (f)--(h);\draw[->,thick,or] (f)--(c);
   \draw[->,thick,or] (h)--(b);\draw[->,thick,or] (c)--(b);
  \end{tikzpicture}
  \end{center}
  \vspace{3mm}\centering{\color{rose}Observer, expérimenter, comprendre.}
\end{frame}

\end{document}
